% Generated draft: requires mathematical and source review.
\documentclass[11pt]{article}
\usepackage[margin=25mm]{geometry}
\usepackage{fontspec}
\setmainfont{Latin Modern Roman}
\usepackage{amsmath,amssymb,mathtools,unicode-math}
\setmathfont{Latin Modern Math}
\usepackage{enumitem,needspace,xurl,hyperref}
\setlength{\parindent}{0pt}
\setlength{\parskip}{5pt}
\setlength{\emergencystretch}{3em}
\newcommand{\N}{\mathbb{N}}
\newcommand{\Z}{\mathbb{Z}}
\newcommand{\Q}{\mathbb{Q}}
\newcommand{\R}{\mathbb{R}}
\newcommand{\C}{\mathbb{C}}
\newcommand{\F}{\mathbb{F}}
\newcommand{\A}{\mathbb{A}}
\newcommand{\PP}{\mathbb{P}}
\newcommand{\E}{\mathbb{E}}
\newcommand{\eps}{\varepsilon}
\newcommand{\dd}{\,\mathrm{d}}
\begin{document}
\section*{UnsolvedMath ID: 2302065}
\textbf{Source problem code:} catalogue entry
\subsection{Definitions and mathematical statement}
Write $f^{(0)}=f$. For entire functions $f,g$ whose indicated denominators are not identically zero, the condition
\[
\frac{f^{(n)}}{g^{(n)}}\text{ extends to an entire function with no zeros}
\]
means that $f^{(n)}$ and $g^{(n)}$ have exactly the same zeros with the same multiplicities. The source does not specify how to interpret a quotient when one or both derivatives vanish identically; that case is not assigned an invented $0/0$ convention.

Hinkkanen's question asks whether there is one integer $k\ge2$ such that for every admissible pair of entire functions $f,g$, if this quotient condition holds for every $n=0,1,\ldots,k$, then either $f/g$ is a nonzero constant or $(f,g)$ has one of the forms
\[
f(z)=e^{az+b},\qquad g(z)=e^{cz+d},
\]
\[
f(z)=A(e^{az}-b),\qquad g(z)=B(e^{-az}-b^{-1}).
\]
The parameters are complex. In the second family $b\ne0$ is necessary for the printed expression; all parameters must make the functions and quotient hypotheses meaningful, so zero prefactors are excluded in an admissible nonzero pair. The source gives no further parameter restrictions. The same universal $k$ must work for all pairs, not depend on $f$ and $g$; the question does not prescribe its value.

The source says the proposer proved, unpublished, that $k=2$ suffices in certain cases, including when $f$ and $g$ have finite order. It also gives
\[
f(z)=(e^{2z}-1)\exp(-i e^z),\qquad
g(z)=(1-e^{-2z})\exp(i e^{-z}),
\]
as showing that $k=3$ is sometimes needed. These are supplied claims about the limitation of $k=2$, not a verified proof that $k=3$ suffices generally. Finite order here has the usual entire-function growth meaning.

A similar meromorphic-function question is also proposed, with the additional printed family
\[
f(z)=A(e^{h(z)}-1)^{-1},\qquad
g(z)=B(1-e^{-h(z)}),
\]
for a nonconstant entire $h$. The second expression is retained exactly without adding a reciprocal. The source does not specify what replaces the entire nonvanishing quotient condition in this meromorphic variant, including treatment of poles and common divisors. Thus it is not made into an unqualified extension of the precise entire-function implication.

The supplied triage labels the optimal-$k$ question open as of the 2018 edition and reports no recent resolution. This is an attributed status assessment, not an independently checked current theorem or proof.
\subsection{Short English statement}
Is there a universal finite number of derivatives whose matching zero divisors determine an entire function up to a scalar and the two stated exceptional families? The source reports unpublished finite-order results for two derivatives and an example requiring at least three in general. It also suggests a meromorphic analogue, but leaves its divisor condition unspecified.
\subsection{Sources}
\begin{itemize}
\item Local source registry: \url{data/unsolvedmath/problems.json}, permanent UnsolvedMath ID $2302065$, source code catalogue entry, including the supplied OPEN-TRIAGE assessment. Dataset attribution: Ulam AI and UnsolvedMath Contributors; original-source attribution and rights are retained.
\item Hayman, \emph{Research Problems in Function Theory}, source-list date 2018, Problem 2.65, attributed to A. Hinkkanen: \url{https://arxiv.org/abs/1809.07200}. The metadata reports source-TeX extraction and July 29, 2026 access, with current-status and rights review flags.
\item Hinkkanen's claimed $k=2$ special cases, including finite order, are explicitly described as unpublished in the statement. The displayed example and meromorphic variant have no separate full reference in the record. No linked source was inspected for this draft.
\end{itemize}
\end{document}
